% !TeX root = closed-systems-from-comparison-completeness.tex
% ============================================================
\chapter{\texorpdfstring{Intrinsic \(2\)-Skeleton Obstruction on the Stone Limit}{Intrinsic 2-Skeleton Obstruction on the Stone Limit}}
\label{ch:interface}
% ============================================================
\section{Introduction}
Under the standing principle of closed-world admissibility
(\cref{sp:closed-world}), this chapter develops the inverse-limit
interface consequences of \cref{sp:compactness,sp:transport}.
Taking the triangle boundary and curvature realization from
\cref{ch:triangles,ch:christoffel} as fixed input, it proves that the
first stable intrinsic transport carrier is the quadratic carrier,
stabilizes that carrier on the Stone limit, and then relates nonzero
stabilized square classes to nonzero realized curvature for
\cref{ch:causality,ch:hilbert,ch:einstein} under faithful smooth
realization on this already-forced carrier.
\Cref{ch:triangles,ch:christoffel} established two structural facts
(\cref{thm:triangle,cor:first-visible-comm,thm:christoffel-main}).
First, the first finite witness of morphism-level transport defect is
triangular. Triangle relations therefore define the intrinsic
\(2\)-skeletal boundary of the transport system.
Second, the first visible graded carrier of intrinsic transport
noncommutativity is the quadratic quotient
\[
	F^2/F^3,
\]
and, under smooth realization, its smooth infinitesimal realization is
Riemann curvature.
The purpose of the present chapter is to connect these two facts
intrinsically, before any smooth realization is imposed.
The mechanism is algebraic.
At each finite stage, the isotropy group carries the augmentation
filtration
\[
	F^1K \supseteq F^2K \supseteq F^3K \supseteq \cdots,
\]
and commutator squares automatically land in the quadratic layer
\(F^2K\). Their images in
\[
	F^2K/F^3K
\]
define canonical degree--\(2\) transport classes.
The main point is that nontrivial triangle obstruction determines the
quadratic carrier, that this determination is forced rather than
postulated, that compatible quadratic square classes are preserved
under refinement and hence stabilize on the Stone limit, and that under
smooth realization faithful on this carrier a nonzero stabilized square
class is equivalent to nonzero realized curvature.
The chapter is organized as follows.
We first make the refinement-compatible transport system explicit at the
level of based loop groups and isotropy groups.
We then define the projective triangle obstruction.
Next we establish quadratic cancellation and define the degree--\(2\)
defect classes.
We then isolate the passage from triangle obstruction to the quadratic
carrier.
After that we prove that triangle-orientation reversal descends to an
intrinsic involution on the stabilized quadratic carrier.
Finally we establish, under smooth realization faithful on the forced
quadratic carrier, that a nonzero stabilized square class is equivalent
to nonzero realized curvature.
Thus the index ``\(2\)'' in the interface is not a hypothesis: it is
the theorem proved by the triangle minimality, commutator-filtration,
and refinement-stabilization chain.  The remaining faithfulness
condition belongs only to the realization map; it says that smooth
realization does not collapse the already forced degree--\(2\) carrier.
% ============================================================
\section{Refinement-compatible transport}
\label{sec:interface-refinement}
Let
\[
	\B_1\subseteq \B_2\subseteq \cdots
\]
be the Boolean refinement tower, and write
\[
	S_k:=\UF(\B_k),
	\qquad
	S_\infty \cong \varprojlim_k S_k
\]
for the associated Stone spaces.
For each stage \(k\), let
\[
	\mathcal N_k=(V_k,E_k)
\]
be a finite directed transport network equipped with a transport functor
\[
	\Hol_k:\Path^\pm(\mathcal N_k)\longrightarrow \mathsf C_k.
\]
Fix a basepoint \(p_k\in V_k\), and define
\[
	\Omega_{p_k}(\mathcal N_k)
	:=
	\Hom_{\Path^\pm(\mathcal N_k)}(p_k,p_k),
	\qquad
	K_{p_k}^{(k)}
	:=
	\Aut_{\mathsf C_k}(p_k).
\]
The restriction of \(\Hol_k\) to based loops is the defect map
\[
	\delta_{p_k}^{(k)}:
	\Omega_{p_k}(\mathcal N_k)\longrightarrow K_{p_k}^{(k)}.
\]
Assume the refinement package provides bonding functors
\[
	R_{k+1,k}:\Path^\pm(\mathcal N_{k+1})\to \Path^\pm(\mathcal N_k),
	\qquad
	Q_{k+1,k}:\mathsf C_{k+1}\to \mathsf C_k,
\]
such that
\begin{equation}\label{eq:interface-compat}
	Q_{k+1,k}\circ \Hol_{k+1}
	=
	\Hol_k\circ R_{k+1,k}.
\end{equation}
\begin{lemma}[Bonding maps on based loops and isotropy]
	\label{lem:interface-bonding}
	The refinement functors induce homomorphisms
	\[
		R_{k+1,k}^{\Omega}:
		\Omega_{p_{k+1}}(\mathcal N_{k+1})\longrightarrow
		\Omega_{p_k}(\mathcal N_k),
		\qquad
		Q_{k+1,k}^{K}:
		K_{p_{k+1}}^{(k+1)}\longrightarrow
		K_{p_k}^{(k)},
	\]
	and these satisfy
	\[
		Q_{k+1,k}^{K}\circ \delta_{p_{k+1}}^{(k+1)}
		=
		\delta_{p_k}^{(k)}\circ R_{k+1,k}^{\Omega}.
	\]
\end{lemma}
\begin{proof}
	Since \(R_{k+1,k}\) is a functor, it sends loops at \(p_{k+1}\) to loops
	at \(p_k\), hence restricts to a homomorphism
	\[
		R_{k+1,k}^{\Omega}:
		\Omega_{p_{k+1}}(\mathcal N_{k+1})\to \Omega_{p_k}(\mathcal N_k).
	\]
	Likewise \(Q_{k+1,k}\) restricts on isotropy groups to
	\[
		Q_{k+1,k}^{K}:
		\Aut_{\mathsf C_{k+1}}(p_{k+1})\to \Aut_{\mathsf C_k}(p_k).
	\]
	Now restrict \cref{eq:interface-compat} to
	\(\Omega_{p_{k+1}}(\mathcal N_{k+1})\). This yields
	\[
		Q_{k+1,k}^{K}\circ \delta_{p_{k+1}}^{(k+1)}
		=
		\delta_{p_k}^{(k)}\circ R_{k+1,k}^{\Omega},
	\]
	as claimed.
\end{proof}
% ============================================================
\section{Triangle obstruction}
\label{sec:interface-triangle}
For each stage \(k\), let
\[
	N_{\triangle,p_k}^{(k)}\trianglelefteq \Omega_{p_k}(\mathcal N_k)
\]
denote the normal closure of all based triangle loops.
\begin{definition}[Finite-stage triangle obstruction]
	Define
	\[
		\Obs_\triangle^{(k)}
		:=
		\delta_{p_k}^{(k)}\bigl(N_{\triangle,p_k}^{(k)}\bigr)
		\subseteq
		K_{p_k}^{(k)}.
	\]
\end{definition}
\begin{proposition}[Projective triangle obstruction]
	The groups \(\Obs_\triangle^{(k)}\) form a projective system under
	refinement. Consequently
	\[
		\Obs_\triangle^\infty
		:=
		\varprojlim_k \Obs_\triangle^{(k)}
	\]
	is canonically defined on the inverse-limit locus.
\end{proposition}
\begin{proof}
	Refinement preserves triangle incidence, hence sends a based triangle
	loop at stage \(k+1\) to a based triangle loop at stage \(k\). Since
	\(N_{\triangle,p_k}^{(k)}\) is the normal closure of the based triangle loops at stage
	\(k\), it follows that
	\[
		R_{k+1,k}^{\Omega}\bigl(N_{\triangle,p_{k+1}}^{(k+1)}\bigr)
		\subseteq
		N_{\triangle,p_k}^{(k)}.
	\]
	Applying \cref{lem:interface-bonding} gives
	\[
		Q_{k+1,k}^{K}\Bigl(\delta_{p_{k+1}}^{(k+1)}\bigl(N_{\triangle,p_{k+1}}^{(k+1)}\bigr)\Bigr)
		=
		\delta_{p_k}^{(k)}\Bigl(R_{k+1,k}^{\Omega}\bigl(N_{\triangle,p_{k+1}}^{(k+1)}\bigr)\Bigr)
		\subseteq
		\delta_{p_k}^{(k)}\bigl(N_{\triangle,p_k}^{(k)}\bigr).
	\]
	That is,
	\[
		Q_{k+1,k}^{K}\bigl(\Obs_\triangle^{(k+1)}\bigr)
		\subseteq
		\Obs_\triangle^{(k)}.
	\]
	Hence the triangle obstruction groups form a projective system.
\end{proof}
\begin{remark}[Minimal \(2\)-skeletal character]
	Triangle loops are the minimal \(2\)-cell boundaries in the intrinsic
	transport complex.
	Accordingly \(\Obs_\triangle^{(k)}\) is the first obstruction that is
	genuinely \(2\)-skeletal rather than merely path-theoretic.
\end{remark}
% ============================================================
\section{Augmentation filtration and quadratic cancellation}
\label{sec:interface-filtration}
Let \(K\) be a group and
\[
	\varepsilon:\mathbb Z[K]\to \mathbb Z
\]
its augmentation homomorphism. Write
\[
	I:=\ker(\varepsilon).
\]
\begin{definition}[Intrinsic augmentation filtration]
	For \(m\ge 1\), define
	\[
		F^mK
		:=
		\{g\in K: g-1\in I^m\}.
	\]
\end{definition}
\begin{lemma}[Commutators raise augmentation order]
	\label{lem:interface-comm-raises}
	For all integers \(r,s\ge 1\),
	\[
		[F^rK,F^sK]\subseteq F^{r+s}K.
	\]
\end{lemma}
\begin{proof}
	Let \(a\in F^rK\) and \(b\in F^sK\). Then
	\[
		a-1\in I^r,
		\qquad
		b-1\in I^s.
	\]
	Write
	\[
		u:=a-1\in I^r,
		\qquad
		v:=b-1\in I^s.
	\]
	Then
	\[
		a=1+u,
		\qquad
		b=1+v.
	\]
	In the \(I\)-adic completion one has convergent geometric-series
	expansions
	\[
		a^{-1}=1-u+u^2-u^3+\cdots,
		\qquad
		b^{-1}=1-v+v^2-v^3+\cdots.
	\]
	Hence
	\[
		a^{-1}-1\in I^r,
		\qquad
		b^{-1}-1\in I^s.
	\]
	Now compute
	\[
		[a,b]-1=aba^{-1}b^{-1}-1.
	\]
	Modulo \(I^{r+s}\), every nonconstant term in the product
	\[
		(1+u)(1+v)\bigl(1+(a^{-1}-1)\bigr)\bigl(1+(b^{-1}-1)\bigr)
	\]
	contains at least one factor from \(I^r\) and at least one factor from
	\(I^s\), or else lies in a strictly higher power of \(I\). Hence
	\[
		[a,b]-1\in I^{r+s}.
	\]
	Therefore
	\[
		[a,b]\in F^{r+s}K.
	\]
\end{proof}
Now fix a stage \(k\) and a basepoint \(p_k\).
For based loops \(e_i,e_j\in \Omega_{p_k}(\mathcal N_k)\), define the
commutator square
\[
	\square_{ij}:=e_ie_je_i^{-1}e_j^{-1}.
\]
\begin{proposition}[Quadratic cancellation]
	For every stage \(k\) and every commutator square \(\square_{ij}\),
	\[
		\delta_{p_k}^{(k)}(\square_{ij})\in F^2K_{p_k}^{(k)}.
	\]
\end{proposition}
\begin{proof}
	Since \(\delta_{p_k}^{(k)}\) is a homomorphism,
	\[
		\delta_{p_k}^{(k)}(\square_{ij})
		=
		[\delta_{p_k}^{(k)}(e_i),\delta_{p_k}^{(k)}(e_j)].
	\]
	Every group element lies in \(F^1K\). Hence
	\cref{lem:interface-comm-raises} with \(r=s=1\) gives
	\[
		\delta_{p_k}^{(k)}(\square_{ij})\in F^2K_{p_k}^{(k)}.
	\]
\end{proof}
\begin{lemma}[Primitive two-cell Magnus term]
\label{lem:primitive-two-cell-magnus}
Let \(a,b\in K\) be two transport germs whose degree--\(1\) classes
\[
\bar a,\bar b\in F^1K/F^2K
\]
are independent.  Then the commutator square
\[
[a,b]=aba^{-1}b^{-1}
\]
has nonzero quadratic class
\[
[a,b]\bmod F^3K\in F^2K/F^3K.
\]
More precisely, in the augmentation ideal one has
\[
[a,b]-1\equiv (a-1)(b-1)-(b-1)(a-1) \pmod {I^3},
\]
so the class of \([a,b]\) is the exterior product
\(\bar a\wedge\bar b\).
\end{lemma}
\begin{proof}
Write \(u=a-1\) and \(v=b-1\).  Modulo \(I^3\),
\[
	a^{-1}\equiv 1-u+u^2,
	\qquad
	b^{-1}\equiv 1-v+v^2.
\]
Multiplying
\[
	(1+u)(1+v)(1-u+u^2)(1-v+v^2)
\]
and discarding terms in \(I^3\) gives
\[
	[a,b]\equiv 1+uv-vu \pmod {I^3}.
\]
Thus
\[
	[a,b]-1\equiv uv-vu \pmod {I^3}.
\]
The antisymmetric tensor \(uv-vu\) is the image of
\(\bar a\wedge\bar b\).  Since \(\bar a\) and \(\bar b\) are independent,
this exterior product is nonzero.  Hence \([a,b]\notin F^3K\), while
\cref{lem:interface-comm-raises} gives \([a,b]\in F^2K\).
\end{proof}
\begin{lemma}[Minimal triangle witnesses have nonzero area class]
\label{lem:minimal-triangle-area}
A primitive triangle witness in the strict triangle regime determines two
first-order edge germs with nonzero exterior product in \(F^1K/F^2K\).
\end{lemma}
\begin{proof}
By \cref{thm:triangle} and \cref{cor:bridge-minimal}, every two-edge
restriction of a primitive triangle witness is coherent, while the full
three-edge subsystem is not.  If the degree--\(1\) edge germs of the
triangle spanned a rank--\(0\) or rank--\(1\) subgroup, then the full
triangle would factor through the same one-dimensional first-order
transport line as its two-edge restrictions.  In that case pairwise
coherence would determine the third edge and would extend to a coherent
vertex gauge on the full triangle, contradicting the strict minimality of
\cref{thm:triangle}.  Hence the primitive triangle witness has a rank--\(2\)
first-order span.  Choosing two independent edge germs \(a,b\) in that span
gives \(\bar a\wedge\bar b\neq0\).
\end{proof}
% ============================================================
\section{Quadratic defect classes}
\label{sec:interface-quadratic}
\begin{definition}[Quadratic defect class]
	For each stage \(k\), define
	\[
		\mathcal A_{ij}^{(k)}(p_k)
		:=
		\bigl[\delta_{p_k}^{(k)}(\square_{ij})\bigr]
		\in
		F^2K_{p_k}^{(k)}/F^3K_{p_k}^{(k)}.
	\]
\end{definition}
\begin{proposition}[Alternation and graded bilinearity]
	The assignment
	\[
		(i,j)\longmapsto \mathcal A_{ij}^{(k)}(p_k)
	\]
	is alternating, and commutator induces a bilinear map
	\[
		\bigl(F^1K_{p_k}^{(k)}/F^2K_{p_k}^{(k)}\bigr)
		\times
		\bigl(F^1K_{p_k}^{(k)}/F^2K_{p_k}^{(k)}\bigr)
		\longrightarrow
		F^2K_{p_k}^{(k)}/F^3K_{p_k}^{(k)}.
	\]
\end{proposition}
\begin{proof}
	By \cref{lem:interface-comm-raises}, the commutator of two degree--\(1\)
	elements lands in degree \(2\). Therefore the commutator descends to a
	map
	\[
		\bigl(F^1K/F^2K\bigr)\times\bigl(F^1K/F^2K\bigr)\to F^2K/F^3K.
	\]
	To prove bilinearity, let \(a,a',b\in F^1K\). The standard commutator
	identity gives
	\[
		[aa',b]=[a,b]\,[a,b,a']\,[a',b],
	\]
	where \([a,b,a']=[[a,b],a']\). Since \([a,b]\in F^2K\), one has
	\[
		[a,b,a']\in [F^2K,F^1K]\subseteq F^3K
	\]
	by \cref{lem:interface-comm-raises}. Hence modulo \(F^3K\),
	\[
		[aa',b]\equiv [a,b][a',b].
	\]
	Similarly,
	\[
		[a,bb']\equiv [a,b][a,b'] \pmod{F^3K}.
	\]
	Thus the induced map is bilinear on the quotients.
	For alternation, one has
	\[
		[a,a]=e
	\]
	for every \(a\), so the induced class of \([a,a]\) is zero in
	\(F^2K/F^3K\). Also,
	\[
		[b,a]=[a,b]^{-1},
	\]
	so in additive notation on the abelian quotient \(F^2K/F^3K\),
	\[
		\mathcal A_{ji}^{(k)}(p_k)=-\mathcal A_{ij}^{(k)}(p_k).
	\]
	Therefore the assignment is alternating.
\end{proof}
% ============================================================
\section{\texorpdfstring{Persistence of triangle obstruction in degree \(2\)}{Persistence of triangle obstruction in degree 2}}
\label{sec:interface-persistence}
The next step is the intrinsic passage from triangle obstruction to a
nonzero quadratic carrier.
\begin{lemma}[Primitive triangle defects survive quadratically]
	\label{lem:quadratic-nondegenerate}
	Let
	\[
		g\neq e
	\]
	be a primitive transport defect arising from a triangle witness.
	Then there exists a refinement stage \(\ell\) such that
	\[
		g\notin F^3K_{p_\ell}^{(\ell)}.
	\]
\end{lemma}
\begin{proof}
	By \cref{lem:minimal-triangle-area}, the primitive triangle witness
	determines two first-order edge germs \(a,b\) with
	\[
		\bar a\wedge\bar b\neq0.
	\]
	The comparison of the two opposite orderings of these germs is the
	associated commutator square.  By
	\cref{lem:primitive-two-cell-magnus}, that square has nonzero image in
	\[
		F^2K/F^3K.
	\]
	Refinement compatibility carries this primitive square class to the
	corresponding square class at sufficiently fine stages.  Therefore the
	refined primitive defect cannot lie in \(F^3\) at every stage.  Hence
	there exists \(\ell\) with
	\[
		g\notin F^3K_{p_\ell}^{(\ell)}.
	\]
\end{proof}
\begin{lemma}[Persistence of nontrivial triangle defect in \(F^2/F^3\)]
	\label{lem:interface-persistence}
	If
	\[
		\Obs_\triangle^{(k)}\neq \{e\},
	\]
	then, after passing if necessary to a sufficiently fine refinement stage
	\(\ell\ge k\), the image of \(\Obs_\triangle^{(k)}\) in
	\[
		F^2K_{p_\ell}^{(\ell)}/F^3K_{p_\ell}^{(\ell)}
	\]
	is nonzero.
\end{lemma}
\begin{proof}
	The group \(\Obs_\triangle^{(k)}\) is generated by the images of based
	primitive triangle loops and their conjugates.  If every primitive
	triangle defect had trivial image, the generated obstruction group would
	be trivial.  Since \(\Obs_\triangle^{(k)}\neq \{e\}\), at least one
	primitive triangle witness has nontrivial defect.

	By refinement compatibility, that primitive defect propagates to finer
	stages.  By \cref{lem:quadratic-nondegenerate}, there exists a refinement
	stage \(\ell\ge k\) at which its refined image does not lie in
	\[
		F^3K_{p_\ell}^{(\ell)}.
	\]
	On the other hand, the witness is a boundary-supported \(2\)-cell
	defect, so its degree--\(1\) image vanishes by
	\cref{cor:triangle-degree1-invisible} in \cref{ch:christoffel}.  Its
	first possible nonzero image is therefore in the quadratic range.  Hence
	its class in
	\[
		F^2K_{p_\ell}^{(\ell)}/F^3K_{p_\ell}^{(\ell)}
	\]
	is nonzero, and so the image of \(\Obs_\triangle^{(k)}\) in this quotient
	is nonzero.
\end{proof}
% ============================================================
\section{The relevant quadratic carrier}
\label{sec:interface-generation}
The next object is the intrinsic degree--\(2\) carrier generated by
quadratic square classes coming from the triangle-induced transport
sector.
\begin{definition}[Finite-stage quadratic carrier]
	\label{def:interface-Kk}
	Let
	\[
		\mathcal K_k
		\subseteq
		F^2K_{p_k}^{(k)}/F^3K_{p_k}^{(k)}
	\]
	denote the subgroup generated by the classes
	\[
		\mathcal A_{ij}^{(k)}(p_k)
	\]
	arising from commutator squares in the triangle-closed transport
	regime.
\end{definition}
\begin{lemma}[Quadratic carrier is square-generated]
	For each stage \(k\), the subgroup
	\[
		\mathcal K_k
		\subseteq
		F^2K_{p_k}^{(k)}/F^3K_{p_k}^{(k)}
	\]
	defined in \cref{def:interface-Kk} is generated by the classes
	\[
		\mathcal A_{ij}^{(k)}(p_k).
	\]
\end{lemma}
\begin{proof}
	This is immediate from the definition of \(\mathcal K_k\).
\end{proof}
\begin{remark}[What is retained from triangle defect]
	The carrier \(\mathcal K_k\) is the intrinsic degree--\(2\) residue of
	the \(2\)-skeletal transport boundary.
	It is exactly the portion of triangle obstruction that remains visible
	after quotienting away cubic and higher transport error.
\end{remark}
% ============================================================
\section{Quadratic carrier determined by triangle obstruction}
\label{sec:interface-carrier}
\begin{theorem}[Triangle obstruction determines a quadratic carrier]
	\label{thm:interface-quad-carrier}
	If
	\[
		\Obs_\triangle^{(k)}\neq \{e\},
	\]
	then after refinement there exist \(\ell\ge k\) and indices
	\((i,j)\) such that
	\[
		\mathcal A_{ij}^{(\ell)}(p_\ell)\neq 0.
	\]
\end{theorem}
\begin{proof}
	By \cref{lem:interface-persistence}, after refinement to some
	\(\ell\ge k\), the image of \(\Obs_\triangle^{(k)}\) in
	\[
		F^2K_{p_\ell}^{(\ell)}/F^3K_{p_\ell}^{(\ell)}
	\]
	is nonzero.
	By definition, the relevant degree--\(2\) carrier is the subgroup
	\(\mathcal K_\ell\) generated by the square classes
	\[
		\mathcal A_{ij}^{(\ell)}(p_\ell).
	\]
	Therefore at least one such class must be nonzero.
\end{proof}
% ============================================================
\section{Universal property of the quadratic carrier}
\label{sec:interface-universal}
\Cref{thm:interface-quad-carrier} shows that the degree--\(2\) quotient is the first
place where intrinsic \(2\)-skeletal obstruction becomes visible.
The next proposition isolates its universal property.
\begin{proposition}[Universal property of the quadratic carrier]
	\label{prop:interface-universal-quadratic}
	Let
	\[
		\pi_{2,3}:F^2K_p\to F^2K_p/F^3K_p
	\]
	be the quotient map.
	Then \(\pi_{2,3}\) is initial among group morphisms from \(F^2K_p\) to abelian
	groups which annihilate \(F^3K_p\).
	Equivalently, if
	\[
		\phi:F^2K_p\to A
	\]
	is a group morphism into an abelian group \(A\) with
	\[
		F^3K_p\subseteq \ker(\phi),
	\]
	then there exists a unique morphism
	\[
		\widetilde\phi:F^2K_p/F^3K_p\to A
	\]
	such that
	\[
		\phi=\widetilde\phi\circ \pi_{2,3}.
	\]
\end{proposition}
\begin{proof}
	This is the universal property of the quotient group
	\(F^2K_p/F^3K_p\).
\end{proof}
\begin{remark}
	\Cref{prop:interface-universal-quadratic} gives the precise sense in
	which \(F^2/F^3\) is the canonical carrier of quadratic transport
	defect.
	Any detector defined on the quadratic layer and insensitive to cubic
	error factors uniquely factors through this quotient.
\end{remark}
% ============================================================
\section{Stabilization on the Stone limit}
\label{sec:interface-stabilization}
\begin{definition}[Stabilized quadratic carrier]
	Define
	\[
		\mathcal K_\infty
		:=
		\varprojlim_k \mathcal K_k.
	\]
\end{definition}
\begin{theorem}[Stabilized quadratic defect]
	\label{thm:interface-stabilized}
	Every compatible family of finite-stage quadratic square classes defines
	an element of
	\[
		\mathcal K_\infty=\varprojlim_k \mathcal K_k.
	\]
	Conversely, every element of \(\mathcal K_\infty\) is represented by a
	compatible family whose \(k\)-th component lies in \(\mathcal K_k\).
\end{theorem}
\begin{proof}
	By \cref{lem:interface-bonding}, the refinement package supplies
	group homomorphisms
	\(Q_{k+1,k}^{K}:K_{p_{k+1}}^{(k+1)}\to K_{p_k}^{(k)}\). Any group
	homomorphism extends linearly to a morphism of group rings preserving
	augmentation ideals, so \(Q_{k+1,k}^{K}\) carries
	\(F^2K_{p_{k+1}}^{(k+1)}\) into \(F^2K_{p_k}^{(k)}\) and
	\(F^3K_{p_{k+1}}^{(k+1)}\) into \(F^3K_{p_k}^{(k)}\). Therefore it
	induces a homomorphism
	\(F^2K_{p_{k+1}}^{(k+1)}/F^3K_{p_{k+1}}^{(k+1)}\to
	F^2K_{p_k}^{(k)}/F^3K_{p_k}^{(k)}\). Because each
	\(\mathcal K_k\) is generated by the square classes
	\(\mathcal A_{ij}^{(k)}(p_k)\) and refinement carries such square
	classes to square classes at the previous stage, these quotient maps
	restrict to bonding maps \(\mathcal K_{k+1}\to\mathcal K_k\). By
	definition of inverse limit, compatible families of finite-stage
	square classes then determine exactly the elements of
	\(\mathcal K_\infty=\varprojlim_k \mathcal K_k\), and every element
	of \(\mathcal K_\infty\) is represented by such a compatible family.
\end{proof}
\begin{definition}[Stabilized square class]
	Whenever a compatible family of finite-stage square classes
	\[
		\bigl(\mathcal A_{ij}^{(k)}(p_k)\bigr)_k
	\]
	is fixed, we write its image in \(\mathcal K_\infty\) as
	\[
		\mathcal A_{ij}(p_\infty).
	\]
\end{definition}
\begin{theorem}[Second-layer stabilization theorem]
	\label{thm:interface-seal}
	In the closed comparison stack, the first stable nontrivial transport
	obstruction in the refinement tower is the stabilized quadratic carrier
	\[
		\mathcal K_\infty\simeq F^2/F^3.
	\]
	Equivalently, a nontrivial triangle obstruction cannot first appear in
	degree \(1\), and it cannot be supported purely in \(F^3\) or in any
	higher filtration layer.  After refinement it has a nonzero image in the
	degree--\(2\) carrier, and compatible such images stabilize on the Stone
	limit.
\end{theorem}
\begin{proof}
	By \cref{thm:triangle} and \cref{cor:bridge-minimal}, the first finite
	morphism-level obstruction is triangular: every two-edge restriction is
	coherent, while the full three-edge subsystem need not be.  Hence no
	degree--\(1\) transport quotient can carry this obstruction.  This is the
	content of \cref{cor:triangle-degree1-invisible} in
	\cref{ch:christoffel}: boundary-supported \(2\)-cell defect has trivial
	image in \(F^1/F^2\).

	By \cref{lem:interface-persistence}, any nontrivial triangle obstruction
	has, after sufficiently fine refinement, a nonzero image in
	\[
		F^2K_{p_\ell}^{(\ell)}/F^3K_{p_\ell}^{(\ell)}.
	\]
	By \cref{thm:interface-quad-carrier}, this nonzero image is represented
	by a nonzero quadratic square class
	\[
		\mathcal A_{ij}^{(\ell)}(p_\ell).
	\]
	Equivalently, the obstruction cannot first appear in degree \(3\) or
	higher: such an appearance would force the degree--\(2\) image to vanish,
	contradicting the nonzero square class just obtained.

	Finally, \cref{thm:interface-stabilized} shows that refinement-compatible
	families of these finite-stage square classes define exactly the inverse
	limit carrier
	\[
		\mathcal K_\infty=\varprojlim_k\mathcal K_k.
	\]
	Thus the first stable nontrivial transport obstruction of the refinement
	tower is neither degree \(1\) nor a hidden higher layer.  It is exactly
	the stabilized quadratic carrier \(\mathcal K_\infty\simeq F^2/F^3\).
\end{proof}
\begin{remark}[What has stabilized]
	The object \(\mathcal K_\infty\) is the intrinsic degree--\(2\)
	transport carrier retained by the full refinement tower.
	It is the first stabilized obstruction visible on the Stone inverse
	limit.
\end{remark}
\begin{definition}[Admissible degree--\(2\) scalar detector]
	\label{def:interface-admissible-scalar}
	Fix a basepoint \(p\). A group morphism
	\[
		\phi:F^2K_p\to \mathbb R
	\]
	is called an admissible degree--\(2\) scalar detector if it is invariant
	under multiplication by elements of \(F^3K_p\), i.e. if
	\[
		\phi(xz)=\phi(x)
		\qquad
		\text{for all }x\in F^2K_p,\ z\in F^3K_p.
	\]
\end{definition}
\begin{theorem}[Kernel annihilation at the quadratic carrier]
	\label{thm:kernel-annihilation-upstream}
	Fix a basepoint \(p\), and let
	\[
		\pi_{2,3}:F^2K_p\to F^2K_p/F^3K_p
	\]
	be the canonical quotient.
	Let
	\[
		\phi:F^2K_p\to \mathbb R
	\]
	be any admissible degree--\(2\) scalar detector.
	Then
	\[
		F^3K_p\subseteq \Ker(\phi),
	\]
	and therefore there exists a unique morphism
	\[
		\overline\phi:F^2K_p/F^3K_p\to\mathbb R
	\]
	such that
	\[
		\phi=\overline\phi\circ \pi_{2,3}.
	\]
\end{theorem}
\begin{proof}
	Let \(z\in F^3K_p\).
	Since \(\phi\) is admissible, \cref{def:interface-admissible-scalar}
	gives
	\[
		\phi(xz)=\phi(x)
		\qquad
		\text{for all }x\in F^2K_p.
	\]
	Taking \(x=e\), where \(e\) is the identity element of \(F^2K_p\),
	yields
	\[
		\phi(z)=\phi(e).
	\]
	Because \(\phi\) is a group morphism into the additive group
	\(\mathbb R\), one has
	\[
		\phi(e)=0.
	\]
	Hence
	\[
		\phi(z)=0
		\qquad
		\text{for all }z\in F^3K_p.
	\]
	Therefore
	\[
		F^3K_p\subseteq \Ker(\phi).
	\]
	Now apply \cref{prop:interface-universal-quadratic}. Since \(\phi\)
	annihilates \(F^3K_p\), there exists a unique morphism
	\[
		\overline\phi:F^2K_p/F^3K_p\to\mathbb R
	\]
	such that
	\[
		\phi=\overline\phi\circ \pi_{2,3}.
	\]
	This is the required factorization.
\end{proof}
% ============================================================
\section{Orientation reversal on the quadratic carrier}
\label{sec:orientation-carrier}
Primitive triangle witnesses admit a natural orientation reversal.
We show that this operation descends to the stabilized quadratic carrier.
\begin{definition}[Primitive orientation reversal]
	Let
	\[
		\omega(p,q,r)
	\]
	denote the defect class associated to the oriented triangle
	\((p,q,r)\).
	Define the reversed class by
	\[
		\mathsf R_\triangle\,\omega(p,q,r):=\omega(p,r,q).
	\]
\end{definition}
\begin{lemma}[Orientation reversal preserves the triangle quotient]
	Let \(w_\triangle(p,q,r)\) denote the triangle boundary word.
	Then reversing orientation sends
	\[
		w_\triangle(p,q,r)\mapsto w_\triangle(p,r,q)=w_\triangle(p,q,r)^{-1}.
	\]
	Hence orientation reversal preserves the normal closure of triangle
	relations.
\end{lemma}
\begin{proof}
	By definition
	\[
		w_\triangle(p,q,r)=[rp][qr][pq].
	\]
	Reversing orientation exchanges the last two vertices and yields the
	inverse word.
	Since the triangle relations are imposed through their normal closure,
	the quotient group is invariant under inversion.
\end{proof}
\begin{theorem}[Orientation involution on the quadratic carrier]
	Primitive triangle reversal descends to a well-defined involution
	\[
		\mathsf R_{\mathrm{or}}:\mathcal K_\infty\longrightarrow \mathcal K_\infty.
	\]
	Moreover
	\[
		\mathsf R_{\mathrm{or}}^2=I.
	\]
\end{theorem}
\begin{proof}
	At each finite stage, the quotient
	\[
		F^2K_{p_k}^{(k)}/F^3K_{p_k}^{(k)}
	\]
	is abelian, so inversion is a well-defined automorphism.
	On square classes it exchanges
	\[
		\mathcal A_{ij}^{(k)}(p_k)
		\longleftrightarrow
		\mathcal A_{ji}^{(k)}(p_k)
		=
		-\mathcal A_{ij}^{(k)}(p_k).
	\]
	By refinement compatibility, these involutions are compatible from stage
	to stage and therefore descend to the inverse limit
	\(\mathcal K_\infty\).
	Applying inversion twice yields the identity, so
	\[
		\mathsf R_{\mathrm{or}}^2=I.
	\]
\end{proof}
\begin{remark}
	The involution \(\mathsf R_{\mathrm{or}}\) exchanges the two orientations of every primitive
	triangle witness.
	Later chapters may use the corresponding \(\pm1\)-eigenspace
	decomposition to organize the defect sector.
\end{remark}
% ============================================================
\section{Smooth realization}
\label{sec:interface-smooth}
Assume now a smooth realization in which square-loop transport admits a
small-loop parameter \(\varepsilon\), and suppose parallel transport
around a coordinate square satisfies
\[
	\Hol(\square_{ij}(\varepsilon))
	=
	\id
	+
	\varepsilon^2 R(\partial_i,\partial_j)
	+
	O(\varepsilon^3).
\]
Fix a limit basepoint \(p_\infty\in S_\infty\), and set
\[
	V:=T_{\iota(p_\infty)}M.
\]
\begin{definition}[Second-jet comparison map]
	\label{def:interface-J2}
	Define
	\[
		\Jtwo:\mathcal K_\infty\to \End(V)
	\]
	to be the map sending the quadratic class of an intrinsic square defect
	to the coefficient of \(\varepsilon^2\) in the realized holonomy
	expansion.
\end{definition}
\begin{lemma}[Well-definedness of the second-jet map]
	The map \(\Jtwo\) is well-defined on \(\mathcal K_\infty\).
\end{lemma}
\begin{proof}
	By \cref{lem:filtration-holonomy-order} in \cref{ch:christoffel}, the
	quotient
	\[
		F^2/F^3
	\]
	records exactly the quadratic holonomy coefficient, while terms of order
	\(3\) and higher are killed by passage to the quotient by \(F^3\).
	Hence two representatives of the same stabilized quadratic class have
	the same \(\varepsilon^2\)-coefficient under smooth realization.
	Therefore \(\Jtwo\) is well-defined.
\end{proof}
\begin{theorem}[Curvature identification]
	\label{thm:interface-curvature}
	For the stabilized quadratic class,
	\[
		\Jtwo\bigl(\mathcal A_{ij}(p_\infty)\bigr)
		=
		R(\partial_i,\partial_j).
	\]
\end{theorem}
\begin{proof}
	By definition,
	\[
		\mathcal A_{ij}(p_\infty)
	\]
	is the degree--\(2\) class of the intrinsic commutator defect of the
	square loop.
	By \cref{def:interface-J2}, \(\Jtwo\) extracts the coefficient of
	\(\varepsilon^2\) in the realized holonomy expansion.
	By \cref{cor:infdef} in \cref{ch:christoffel}, the coefficient of
	\(\varepsilon^2\) in the holonomy expansion of a small square is
	exactly
	\[
		R(\partial_i,\partial_j).
	\]
	Therefore
	\[
		\Jtwo\bigl(\mathcal A_{ij}(p_\infty)\bigr)
		=
		R(\partial_i,\partial_j).
	\]
\end{proof}
% ============================================================
\section{Interface theorem}
\label{sec:interface-main}
\begin{theorem}[Intrinsic--smooth interface]
	\label{thm:interface-main}
	Under the standing principle \cref{sp:closed-world}, the refinement
	tower stabilizes its first nontrivial transport obstruction at the
	degree--\(2\) carrier
	\[
		\mathcal K_\infty\simeq F^2/F^3.
	\]
	Assume only realization faithfulness on this forced carrier, i.e. that
	the comparison map
	\[
		\Jtwo:\mathcal K_\infty\to \End(V)
	\]
	is injective.
	Then the following are equivalent:
	\[
		\exists\,\mathcal A_{ij}(p_\infty)\neq0,
		\qquad
		R\neq0.
	\]
\end{theorem}
\begin{proof}
	The stabilization assertion is \cref{thm:interface-seal}.  Thus the
	appearance of the second layer is not an additional interface assumption;
	it is the intrinsic output of the refinement tower.  What remains to
	check is the faithful smooth reading of that carrier.

	Suppose that
	\[
		\mathcal A_{ij}(p_\infty)\neq 0.
	\]
	Since \(\Jtwo\) is injective,
	\[
		\Jtwo\bigl(\mathcal A_{ij}(p_\infty)\bigr)\neq 0.
	\]
	By \cref{thm:interface-curvature},
	\[
		\Jtwo\bigl(\mathcal A_{ij}(p_\infty)\bigr)
		=
		R(\partial_i,\partial_j),
	\]
	hence
	\[
		R\neq 0.
	\]
	Conversely, if
	\[
		R\neq 0,
	\]
	then for some indices \(i,j\),
	\[
		R(\partial_i,\partial_j)\neq 0.
	\]
	By \cref{thm:interface-curvature},
	\[
		\Jtwo\bigl(\mathcal A_{ij}(p_\infty)\bigr)=R(\partial_i,\partial_j)\neq 0.
	\]
	Since the zero class has zero second-jet coefficient by
	\cref{def:interface-J2}, this implies
	\[
		\mathcal A_{ij}(p_\infty)\neq 0.
	\]
\end{proof}
\begin{remark}[Status of second-layer stabilization]
\label{rem:second-jet-status}
The aggressive question is not whether one may assume the second layer;
the answer is that one may not.  The second layer is forced.  The proof is
the chain
\[
\begin{aligned}
&\text{strict triangle minimality}
\Longrightarrow
\text{no degree--1 residue}
\\
&\Longrightarrow
\text{commutator-square defect in }F^2
\Longrightarrow
\text{nonzero class in }F^2/F^3
\\
&\Longrightarrow
\mathcal K_\infty.
\end{aligned}
\]
The only remaining hypothesis is not that the tower stabilizes at degree
\(2\), but that a chosen compatible smooth realization is faithful on the
already stabilized degree--\(2\) carrier.  Under that faithful reading,
\cref{thm:interface-curvature} identifies the realized image with Riemann
curvature, and the GR chain proceeds.
\end{remark}
% ============================================================
\section{Einstein boundary}
\label{sec:interface-boundary}
The present chapter identifies the stabilized quadratic square-class
channel as the forced intrinsic degree--\(2\) transport carrier and,
under faithful smooth realization on that carrier, shows that nonzero
such classes are equivalent to nonzero realized curvature.
No dynamical equations have been assumed or derived.
Macroscopic compatibility laws belong to the subsequent chapter.

\section{Conclusion}
\label{sec:interface-conclusion}
The interface result is now closed in stabilized form: the refinement
tower has a forced degree--\(2\) first obstruction, and under faithful
smooth realization on that carrier the existence of a nonzero stabilized
quadratic square class is equivalent to nonzero realized curvature on the
same limit object.

Accordingly, \cref{ch:causality} combines the stabilized quadratic
carrier, under the faithful-realization criterion that nonzero stabilized
square classes are equivalent to nonzero realized curvature, with the
observer-side filtration inherited from
\cref{ch:flow} and interleaved with the refinement tower in
\cref{ch:stitching} to derive the macroscopic causal and scaling
consequences forced by the closed stack.